% 第一章 引言
\section{引言}

\subsection{研究背景}

数字化浪潮已深刻重塑个人资产的形态。从加密货币、数字钱包到云存储账号、社交平台、加密文件，数字资产在个人财富中的占比持续攀升。据 Chainalysis 2023 年度报告显示，全球加密货币市场规模已突破 15 万亿美元，而 DeFi 生态锁仓价值（TVL）在 2021 年曾达到 1700 亿美元的历史峰值。然而，数字资产在安全性方面存在与传统资产截然不同的挑战：

\begin{enumerate}[label=(\arabic*)]
    \item \textbf{私钥单点风险}：加密货币私钥由持有人独自掌握，一旦持有人去世或失能，资产永久锁定。据估算，仅比特币领域因持有者去世导致的永久性资产损失已超过 2000 亿美元。
    \item \textbf{数字资产孤岛}：云存储、社交平台、邮箱等服务商各有独立的隐私政策，继承人需逐一走繁琐的法律流程才能获取访问权限。
    \item \textbf{胁迫攻击威胁}：持有人可能在生命威胁下被迫转移资产，传统方案缺乏有效的反制手段。
    \item \textbf{传统继承工具失效}：公证遗嘱流程缓慢、成本高昂，且难以覆盖数字资产的时效性和技术复杂性。
\end{enumerate}

\subsection{现有方案及其局限性}

目前，数字资产继承领域已有若干探索性方案，但均存在不同程度的局限性：

\textbf{方案一：Ledger Recover}。Ledger 公司提供的付费订阅服务，将用户私钥加密分片托管给三家机构。其问题在于：第一，该服务为闭源方案，密码学协议不可公开审计；第二，用户需通过 KYC 身份验证，隐私性差；第三，每月 \$9.99 的费用增加了使用门槛；第四，依赖 Ledger 公司的中心化服务，存在单点故障风险。

\textbf{方案二：多签钱包（如 Safe\{Wallet\}）}。基于智能合约的多签方案，设定 m-of-n 签名策略进行资产转移。其优势在于链上透明、去中心化。但局限性同样明显：第一，仅支持 EVM 生态资产，无法覆盖云账号、加密文件等非链上资产；第二，所有交易在链上公开，隐私保护不足；第三，缺乏抗胁迫设计，攻击者可逐个胁迫签名人。

\textbf{方案三：传统公证遗嘱}。通过公证处进行资产继承的法律流程。但数字资产的继承往往涉及密码材料的交割，公证流程无法处理"私钥传递"这一技术环节。且公证流程费用高昂（数千至数万元）、周期长（数月），不满足数字资产继承的时效性需求。

\textbf{方案四：密码管理器继承功能}。如 Bitwarden、1Password 等密码管理器提供了紧急联系人功能。但此类方案存在两个致命缺陷：第一，密码管理器本身构成新的单点故障——若主密码丢失或被破解，所有资产全部泄露；第二，缺乏门槛机制，紧急联系人可直接获取完整密码，存在滥用风险。

\subsection{本文贡献}

针对上述问题，本文提出并实现了\textbf{数字遗产管家}（Digital Legacy Guardian），一个基于门限密码学与同态承诺的可验证数字资产继承协议。本文的主要贡献如下：

\begin{enumerate}[label=(\arabic*)]
    \item \textbf{双多项式份额结构}：在标准 Shamir 秘密共享基础上引入第二个盲因子多项式 Q(x)，使得份额值通过 Pedersen 承诺在公开存储的同时具备同态可验证性。任意第三方可利用 Lagrange 系数在椭圆曲线上直接验证份额的一致性，无需知道份额值本身。

    \item \textbf{零信任同态份额刷新}：提出了基于零和多项式的同态份额刷新协议。在 server \textbf{完全不存储主密钥}的前提下，通过生成 Z(0) = 0 的零和多项式，计算公开增量并利用 Pedersen 承诺的同态加法性质更新承诺。该方案消除了传统方案中"解密→重拆分→重新加密"对主密钥的依赖。

    \item \textbf{抗胁迫胁迫码机制}：设计了计算不可区分的胁迫码方案——每位监护人在正常份额之外额外持有胁迫份额。两者均以 Pedersen 承诺形式存储，攻击者无法区分。提交胁迫份额会暗中触发警报并终止继承流程，在人身威胁场景中提供隐蔽的求救通道。

    \item \textbf{完整系统实现}：基于 Node.js/Express 后端和 React/TypeScript 前端实现了完整的数字资产继承系统，覆盖"创建—分发—验证—刷新—恢复"全生命周期。项目以 MIT 许可证全栈开源，独立开发代码行数为 9177 行（不含第三方库和 node\_modules），具备实际部署能力。
\end{enumerate}

\subsection{论文组织}

本文的后续组织结构如下：第 2 章介绍相关背景知识与密码学基础；第 3 章阐述系统架构与协议设计；第 4 章详细描述技术实现；第 5 章分析安全性并与现有方案对比；第 6 章展示实验评估结果\footnote{需补充实验数据}；第 7 章总结全文并展望未来工作。
